\chapter{1999年上海卷（理）}
\section{填空题}

\begin{problem}
\(\tan\left(\arccos \frac{\sqrt{2}}{2} - \frac{\pi}{6}\right)=\)\fillinblank{}.

\end{problem}

\begin{problem}
函数 \(f(x)=\log_2 x+1\)（\(x\geqslant 4\)）的反函数 \(f^{-1}(x)\) 的定义域是 \fillinblank{}.

\end{problem}

\begin{problem}
在 \(\left( x^3 + \frac{2}{x^2} \right)^5\) 的展开式中，含 \(x^5\) 项的系数为 \fillinblank{}.

\end{problem}

\begin{problem}
在 \(\triangle ABC\) 中，若 \(\angle B=30^\circ\)，\(AB=2\sqrt{3}\)，\(AC=2\)，则 \(\triangle ABC\) 的面积 \(S\) 是 \fillinblank{}.

\end{problem}

\begin{problem}
若平移坐标系，将曲线方程 \(y^2 + 4x - 4y - 4 = 0\) 化为标准方程，则坐标原点应移到点 \(O'\left( \fillinblank{}, \fillinblank{} \right)\).

\end{problem}

\begin{problem}
函数 \(y = 2\sin x \cos x - 2\sin^2 x + 1\) 的最小正周期 \(T=\)\fillinblank{}.

\end{problem}

\begin{problem}
函数 \(y = 2\sin\left( 2x + \frac{\pi}{6} \right)\)（\(x \in \left[ -\pi, 0 \right]\)）的单调递减区间是 \fillinblank{}.

\end{problem}

\begin{problem}
若将向量 \(\bs{a} = \{2, 1\}\) 围绕原点按逆时针方向旋转 \(\frac{\pi}{4}\) 得到向量 \(\bs{b}\)，则向量 \(\bs{b}\) 的坐标为 \fillinblank{}.

\end{problem}

\begin{problem}
极坐标方程 \(5\rho^2 \cos 2\theta + \rho^2 - 24 = 0\) 所表示的曲线焦点的极坐标为 \fillinblank{}.

\end{problem}

\begin{problem}
在等差数列 \(\{a_n\}\) 中，满足 \(3a_4 = 7a_7\)，且 \(a_1 > 0\)，\(S_n\) 是数列 \(\{a_n\}\) 前 \(n\) 项的和. 若 \(S_n\) 取得最大值，则 \(n=\fillinblank{}\).

\end{problem}

\begin{problem}
若以连续掷两次骰子分别得到的点数 \(m\)、\(n\) 作为点 \(P\) 的坐标，则点 \(P\) 落在圆 \(x^2 + y^2 = 16\) 内的概率是 \fillinblank{}.

\end{problem}

\begin{problem}
若四面体各棱的长是 \(1\) 或 \(2\)，且该四面体不是正四面体，则其体积的值是 \fillinblank{}（只需写出一个可能的值）.

\end{problem}

\section{单选题}

\begin{problem}
直线 \(y = \frac{\sqrt{3}}{3}x\) 绕原点按逆时针方向旋转 \(30^\circ\) 后所得直线与圆 \(\left( x - 2 \right)^2 + y^2 = 3\) 的位置关系是……（\quad）

\choices
    {直线过圆心}
    {直线与圆相交，但不过圆心}
    {直线与圆相切}
    {直线与圆没有公共点.}

\end{problem}

\begin{problem}
下列以 \(t\) 为参数的参数方程所表示的曲线中，与方程 \(xy = 1\) 所表示的曲线完全一致的是……（\quad）

\choices
    {\(\begin{cases} x = t^{\frac{1}{2}} \\ y = t^{-\frac{1}{2}} \end{cases}\)}
    {\(\begin{cases} x = |t| \\ y = \frac{1}{|t|} \end{cases}\)}
    {\(\begin{cases} x = \cos t \\ y = \sec t \end{cases}\)}
    {\(\begin{cases} x = \tan t \\ y = \cot t \end{cases}\)}

\end{problem}

\begin{problem}
若 \(a < b < 0\)，则下列结论中正确的是（\quad）

\choices
    {不等式 \(\frac{1}{a} > \frac{1}{b}\) 和 \(\frac{1}{|a|} > \frac{1}{|b|}\) 均不能成立}
    {不等式 \(\frac{1}{a - b} > \frac{1}{a}\) 和 \(\frac{1}{|a|} > \frac{1}{|b|}\) 均不能成立}
    {不等式 \(\frac{1}{a - b} > \frac{1}{a}\) 和 \(\left(a + \frac{1}{b}\right)^2 > \left(b + \frac{1}{a}\right)^2\) 均不能成立}
    {不等式 \(\frac{1}{|a|} > \frac{1}{|b|}\) 和 \(\left(a + \frac{1}{b}\right)^2 >\) \(\left(b + \frac{1}{a}\right)^2\) 均不能成立.}

\end{problem}

\begin{problem}
设有直线 \(m, n\) 和平面 \(\alpha, \beta\)，则在下列命题中，正确的是……（\quad）

\choices
    {若 \(m \parallel n\)，\(m \subset \alpha\)，\(n \subset \beta\)，则 \(\alpha \parallel \beta\)}
    {若 \(m \perp \alpha\)，\(m \perp n\)，\(n \subset \beta\)，则 \(\alpha \parallel \beta\)}
    {若 \(m \parallel n\)，\(n \perp \beta\)，\(m \subset \alpha\)，则 \(\alpha \perp \beta\)}
    {若 \(m \parallel n\)，\(m \perp \alpha\)，\(n \perp \beta\)，则 \(\alpha \perp \beta\).}

\end{problem}

\section{解答题}

\begin{problem}
设集合 \(A = \{x \mid |x - a| < 2\}\)，\(B = \left\{ x \mid \frac{2x - 1}{x + 2} < 1 \right\}\). 若 \(A \subseteq B\)，求实数 \(a\) 的取值范围.

\end{problem}

\begin{problem}
设正数数列 \(\{a_n\}\) 为一等比数列，且 \(a_2 = 4\)，\(a_4 = 16\)，求 \(\displaystyle\lim_{n \to \infty} \frac{\lg a_{n + 1} + \lg a_{n + 2} + \cdots + \lg a_{2n}}{n^2}\).

\end{problem}

\begin{problem}
已知方程 \(x^2 + (4 + \i)x + 4 + a\i = 0\)（\(a \in \R\)）有实根 \(b\)，且 \(z = a + b\i\)，求复数 \(z(1 - c\i)\)（\(c > 0\)）的辐角主值的取值范围.

\end{problem}

\begin{problem}
如图，在四棱锥 \(P-ABCD\) 中，底面 \(ABCD\) 是一直角梯形，\(\angle BAD = 90^\circ\)，\(AD \parallel BC\)，\(AB = BC = a\)，\(AD = 2a\)，且 \(PA \perp\) 底面 \(ABCD\)，\(PD\) 与底面成 \(30^\circ\) 角. 若 \(AE \perp PD\)，\(E\) 为垂足.

\begin{enumerate}
\item 求证：\(BE \perp PD\)；
\item 求异面直线 \(AE\) 与 \(CD\) 所成角的大小（用反三角函数表示）.
\end{enumerate}

\end{problem}

\begin{problem}
平地上有一条水沟，沟沿是两条长 \(100\mathrm{m}\) 的平行线段，沟宽 \(AB\) 为 \(2\mathrm{m}\). 与沟沿垂直的平面与沟的交线是一段抛物线，抛物线顶点为 \(O\)，对称轴与地面垂直，沟深 \(1.5\mathrm{m}\)，沟中水深 \(1\mathrm{m}\).

\begin{enumerate}
\item 求水面宽；
\item 如图所示形状的几何体称为柱体. 已知柱体的体积为底面积乘以高，问沟中的水有多少立方米；
\item 若要把这条水沟改挖（不准填土）成截面为等腰梯形的沟，使沟的底面与地面平行，则改挖后的沟底宽为多少米时，所挖的土最少.
\end{enumerate}
\end{problem}

\begin{problem}
设椭圆 \(C_1\) 的方程为 \(\frac{x^2}{a^2} + \frac{y^2}{b^2} = 1\)（\(a > b > 0\)），曲线 \(C_2\) 的方程为 \(y = \frac{1}{x}\)，且 \(C_1\) 与 \(C_2\) 在第一象限内只有一个公共点 \(P\).

\begin{enumerate}
\item 试用 \(a\) 表示点 \(P\) 的坐标；
\item 设 \(A\)、\(B\) 是椭圆 \(C_1\) 的两个焦点，当 \(a\) 变化时，求 \(\triangle ABP\) 的面积函数 \(S(a)\) 的值域；
\item 记 \(\min\{y_1, y_2, \cdots, y_n\}\) 为 \(y_1, y_2, \cdots, y_n\) 中最小的一个. 设 \(g(a)\) 是以椭圆 \(C_1\) 的半焦距为边长的正方形的面积，试求函数 \(f(a) = \min\{g(a), S(a)\}\) 的表达式.
\end{enumerate}
\end{problem}
